2.2-1 \\
Express the function $n^3/1000 + 100n^2 - 100n + 3$ in terms of $\theta$-notation. \\

2.2-2 \\
Consider sorting n numbers stored in array A[1:n] by first finding the smallest element of A[1:n]
and exchanging it with the element in A[1]. Then find the smalles element of A[2:n], and exchange it with A[2]. 
Then find the smallest element of A[3:n], and exchange it with A[3]. Continue this for the first n-1 elements of A.
Write the code for this algorithm (this is selection sort). What loop invariant does this algorithm maintain? Why does it need
to run for only the first n-1 elements, rantehr than all n? Give the worst-case running time of selection sort in $\theta -notation$. Is the best-case
run time any better? \\

2.2-3 \\
Consider linear search again (see Exercise 2.1-4). How many elements of the input 
array need to be checked on the average, assuming that the element being searched 
for is equally likely to be any element in the array? How about in the worst case? 
Using $\theta -notation$, give the average-case and worst-case run times of linear 
search. Justify your answers. \\

2.2-4 \\
How can you modify any soerting algorithm to have a good best-case running time? 
You can choose from a wide range of algorithm design techniques. Insertion sort 
uses the incremental method: for each element A[i], insert it into its proper place
in the subarray A[1:i], having already sorted the subarray A[1:i-1]. \\
This section examines another design method, known as divide-and-conquer 
which we explore in more detail in Chapter 4. We$’$ll use divide-and-conquer to 
design a sorting algorithm whose worst-case running time is much less than that 
of insertion sort. One advantage of using an algorithm that follows the divide-and- 
conquer method is that analyzing its running time is often straightforward, using 
techniques that we$’$ll explore in Chapter 4. \\


